\chapter{Pricing and Promotions with ML + Optimization}

\section{Learning objectives}
\begin{itemize}
  \item Explain why pricing in e-commerce is both a prediction problem and an optimization problem.
  \item Distinguish demand forecasting, price elasticity estimation, and promotion uplift modeling.
  \item Analyze how business constraints such as inventory, margins, contracts, and competition shape pricing decisions.
  \item Design a pricing workflow that combines forecasting, optimization, and experimentation rather than treating price as a static input.
\end{itemize}

\section{Why pricing is a core AI problem in commerce}
Pricing is one of the most consequential decisions in e-commerce.
It affects demand, conversion, margin, inventory movement, customer perception, and competitive position.
Unlike many prediction tasks, pricing also changes the system it measures.
When a platform changes price, user behavior changes, competitor response may change, and future training data are altered.

This is why pricing cannot be treated as ``just another regression target.''
The practical pricing problem combines:
\begin{itemize}
  \item demand estimation,
  \item elasticity understanding,
  \item business-rule enforcement,
  \item optimization under constraints,
  \item and experimentation to measure causal impact.
\end{itemize}

In enterprise settings, pricing also interacts with ERP-owned price lists, promotions engines, finance controls, and B2B contractual terms.

\section{Pricing objectives}
Different organizations optimize different pricing objectives.
Students should recognize that ``best price'' depends on the business question.

\subsection{Common objectives}
\begin{itemize}
  \item \textbf{Revenue maximization:} increase gross sales volume.
  \item \textbf{Margin or contribution optimization:} emphasize profitability rather than top-line revenue.
  \item \textbf{Inventory balancing:} accelerate slow-moving stock or protect scarce inventory.
  \item \textbf{Customer acquisition or retention:} use pricing to influence long-term account value rather than immediate transaction value.
  \item \textbf{Competitive positioning:} maintain price perception within a market segment.
\end{itemize}

\subsection{Why objective clarity matters}
The same product may warrant different prices depending on whether the business prioritizes:
\begin{itemize}
  \item clearing inventory,
  \item preserving premium brand position,
  \item growing category share,
  \item protecting margin under rising costs.
\end{itemize}

A pricing model trained without explicit objective alignment often produces analytically plausible but operationally poor recommendations.

\section{Price, promotion, and discount structure}
In e-commerce, price is rarely a single number.
The platform may have:
\begin{itemize}
  \item list price,
  \item sale price,
  \item coupon-adjusted price,
  \item customer-specific or B2B negotiated price,
  \item bundle price,
  \item shipping-inclusive effective price.
\end{itemize}

\subsection{Promotions as structured interventions}
Promotions are not merely ``lower prices.''
They can take many forms:
\begin{itemize}
  \item percentage discount,
  \item fixed-amount coupon,
  \item buy-one-get-one structures,
  \item free shipping thresholds,
  \item loyalty-only offers,
  \item category-wide or seller-funded campaigns.
\end{itemize}

Each structure affects user behavior differently.
For example, a coupon may increase conversion among already interested users, while a shipping threshold may increase basket size rather than item-level demand.

\section{Demand forecasting}
Pricing decisions require an estimate of how demand changes over time and under different commercial conditions.

\subsection{What is being forecast}
Forecast targets may include:
\begin{itemize}
  \item unit demand by SKU and day,
  \item category demand,
  \item promotion-period lift,
  \item expected cart conversion,
  \item replenishment-sensitive demand.
\end{itemize}

\subsection{Inputs to demand models}
Typical features include:
\begin{itemize}
  \item historical sales,
  \item current and historical prices,
  \item promotions,
  \item seasonality and holidays,
  \item catalog position and visibility,
  \item marketing traffic,
  \item competitor signals where available,
  \item inventory availability.
\end{itemize}

\subsection{Why forecasting is not enough}
Forecasting predicts what may happen.
Pricing still requires a decision about what the system should do.
A forecast without an optimization layer does not tell the platform which price to choose among feasible options.

\section{Price elasticity}
Elasticity describes how demand changes in response to price changes.
It is one of the most important concepts in pricing analytics.

\subsection{Intuition}
If a small increase in price causes a large reduction in demand, the product is relatively price-sensitive.
If demand changes little, the product is relatively price-insensitive.

\subsection{Why elasticity is difficult to estimate}
In live commerce systems, observed prices are not randomly assigned.
They may reflect prior business policy, inventory conditions, promotions, or manual intervention.
This means historical price-demand relationships are often confounded.

\subsection{Practical considerations}
Elasticity may vary by:
\begin{itemize}
  \item category,
  \item season,
  \item customer segment,
  \item brand strength,
  \item competitive intensity,
  \item stock scarcity,
  \item whether the user arrived organically or through a promotion.
\end{itemize}

Students should therefore avoid assuming a single global elasticity coefficient is sufficient for realistic commerce settings.

\section{Promotions and uplift}
Promotion modeling asks not only ``what happened under promotion?'' but ``what additional effect did the promotion cause?''

\subsection{Observed uplift versus causal uplift}
If conversion rises during a campaign, that increase may reflect:
\begin{itemize}
  \item genuine causal impact from the promotion,
  \item seasonal demand shifts,
  \item traffic-source changes,
  \item selection bias in who received the offer.
\end{itemize}

True uplift modeling attempts to estimate incremental effect rather than raw observed response.

\subsection{Why uplift matters}
Promotions can be expensive.
If the platform discounts heavily for users who would have purchased anyway, the business sacrifices margin without creating incremental value.
Targeted promotion strategies therefore require causal thinking, not just descriptive analytics.

\section{Optimization under business constraints}
After forecasting and elasticity estimation, the pricing problem becomes an optimization problem.

\subsection{Typical constraints}
Real pricing decisions are constrained by:
\begin{itemize}
  \item minimum margin thresholds,
  \item inventory availability,
  \item supplier or brand pricing rules,
  \item MAP or contractual restrictions,
  \item competitive guardrails,
  \item fairness or anti-discrimination rules,
  \item B2B price-list obligations,
  \item operational simplicity limits on how often prices can change.
\end{itemize}

\subsection{Why unconstrained optimization is unrealistic}
A mathematically optimal price from a simple model may be unusable if it:
\begin{itemize}
  \item violates margin policy,
  \item creates stockouts,
  \item damages brand trust,
  \item conflicts with contractual pricing,
  \item or overwhelms operations with constant price updates.
\end{itemize}

Optimization in enterprise commerce therefore means ``best feasible decision,'' not ``best unconstrained number.''

\section{Dynamic pricing}
Dynamic pricing adjusts prices over time in response to changing conditions.
It is widely discussed, but it should not be treated as a universal default.

\subsection{When dynamic pricing can help}
Dynamic pricing is often useful when:
\begin{itemize}
  \item demand fluctuates quickly,
  \item inventory is perishable or time-sensitive,
  \item competitive prices move frequently,
  \item the catalog is large enough that manual pricing is infeasible.
\end{itemize}

\subsection{When dynamic pricing can hurt}
It can also create problems:
\begin{itemize}
  \item customer trust may decline if prices feel unstable or unfair,
  \item operations may struggle with frequent updates,
  \item historical comparison becomes harder,
  \item legal or contractual issues may arise in some markets.
\end{itemize}

The correct architecture therefore distinguishes between categories that justify dynamic treatment and those that require more stable commercial rules.

\section{Competition-aware pricing}
In many categories, prices are influenced by competitor behavior.
Yet competitor-aware pricing is not the same as ``match the lowest price.''

\subsection{Competitive signals}
Useful signals may include:
\begin{itemize}
  \item competitor list price,
  \item promotion intensity,
  \item assortment overlap,
  \item shipping speed comparisons,
  \item seller reputation and availability differences.
\end{itemize}

\subsection{Strategic caution}
A competitor may use a loss leader or temporary campaign that should not be copied directly.
Blind matching can destroy margin without improving long-term position.
The architecture should therefore support competitor monitoring while still respecting internal pricing objectives and constraints.

\section{B2C versus B2B pricing}
Pricing design differs sharply between B2C and B2B commerce.

\subsection{B2C}
B2C pricing usually emphasizes:
\begin{itemize}
  \item high-volume SKU-level decisions,
  \item campaign responsiveness,
  \item public price perception,
  \item consumer conversion behavior.
\end{itemize}

\subsection{B2B}
B2B pricing often involves:
\begin{itemize}
  \item negotiated price lists,
  \item account-specific discounts,
  \item credit terms,
  \item approval workflows,
  \item contract validity windows,
  \item product-assortment restrictions.
\end{itemize}

This means many B2B pricing decisions belong closer to ERP and enterprise sales workflows than to the storefront alone.

\section{Feature design for pricing models}
Pricing models often require features from multiple systems.
Representative feature groups include:
\begin{itemize}
  \item \textbf{Product features:} category, brand, cost, lifecycle stage, substitutes.
  \item \textbf{Demand features:} recent sales velocity, conversion rate, traffic.
  \item \textbf{Promotion features:} campaign type, discount depth, channel.
  \item \textbf{Inventory features:} on-hand stock, days of supply, replenishment lead time.
  \item \textbf{Customer features:} segment, loyalty status, contract class.
  \item \textbf{Market features:} competitor price, holiday periods, macro signals.
\end{itemize}

\subsection{Feature governance}
Pricing is high impact, so feature governance matters.
If cost, inventory, or contract data are delayed or inconsistent, the resulting recommendations may be financially harmful.
This ties pricing directly to the data-foundation and systems-of-record questions from earlier chapters.

\section{Experimentation in pricing}
Offline evaluation is useful, but pricing decisions have strong causal and behavioral effects.
Therefore online experimentation is especially important.

\subsection{Why pricing experiments are sensitive}
Pricing tests differ from many UI experiments because they:
\begin{itemize}
  \item affect revenue directly,
  \item may trigger competitor response,
  \item can influence customer trust,
  \item may have spillover across channels or segments,
  \item often require tighter legal and policy review.
\end{itemize}

\subsection{What to measure}
Useful experiment metrics may include:
\begin{itemize}
  \item conversion rate,
  \item unit sales,
  \item revenue,
  \item contribution margin,
  \item basket size,
  \item repeat purchase,
  \item return rate,
  \item customer support complaints,
  \item price-perception surveys where available.
\end{itemize}

\subsection{Guardrails}
Pricing experiments should include guardrails such as:
\begin{itemize}
  \item margin floor protection,
  \item inventory depletion limits,
  \item customer fairness checks,
  \item B2B account or contractual exclusions.
\end{itemize}

\section{A practical pricing architecture}
A practical pricing system often includes:
\begin{itemize}
  \item demand and elasticity data pipelines,
  \item forecasting models,
  \item optimization logic,
  \item a pricing decision service,
  \item a promotions engine,
  \item ERP integration for governed price lists and approvals,
  \item experimentation and monitoring controls.
\end{itemize}

\subsection{Decision flow}
A simplified pricing flow may be:
\begin{enumerate}
  \item gather product, demand, inventory, and market signals,
  \item forecast demand under candidate price scenarios,
  \item optimize within business constraints,
  \item publish recommended price or promotional action,
  \item validate through workflow or approval gates,
  \item serve price to commerce channels,
  \item measure downstream outcomes.
\end{enumerate}

\subsection{Why architecture matters}
Pricing is not just a model output.
It is a governed operational decision that affects customers, suppliers, finance, and brand perception.
The architecture must therefore support auditability and rollback, not only prediction accuracy.

\section{Common failure modes}
Pricing systems often fail in predictable ways:
\begin{itemize}
  \item forecasting ignores stockouts and mistakes absence of sales for lack of demand,
  \item elasticity estimates are biased by historical policy,
  \item optimization recommends prices that violate margin or contractual constraints,
  \item promotions are evaluated on revenue alone while margin collapses,
  \item competitor signals are stale or incomparable,
  \item different channels display conflicting prices due to synchronization delays.
\end{itemize}

These failures are reminders that pricing is a socio-technical control problem, not just an analytics exercise.

\section{Hands-on lab: forecast plus constrained optimization}
The practical goal of this chapter is to connect prediction and decision-making.

\subsection{Lab scenario}
Students are given historical SKU sales, prices, promotion flags, inventory levels, and seasonality indicators.
They are asked to produce a pricing recommendation for a small catalog segment.

\subsection{Lab tasks}
\begin{enumerate}
  \item Build a baseline demand forecast for a selected product or category.
  \item Estimate how demand changes across a small set of candidate prices.
  \item Define a margin floor and one inventory-related constraint.
  \item Choose a recommended price or promotion and justify it.
  \item Specify how you would validate the decision using an online experiment.
\end{enumerate}

\subsection{Expected learning outcome}
By the end of the lab, students should be able to explain why pricing requires a pipeline of assumptions, constraints, and measurement rather than a single predictive model.

\section{Managerial implications}
Pricing systems can create large business value, but they can also create large business damage when poorly governed.
Organizations should therefore invest in:
\begin{itemize}
  \item clear pricing objectives,
  \item data quality around cost, stock, and catalog attributes,
  \item explicit separation between descriptive analytics and prescriptive decisions,
  \item and controlled experimentation before scaling automated pricing.
\end{itemize}

The most effective pricing systems are not those that change the most prices.
They are those that align predictive intelligence with disciplined commercial control.

\section{Multiple-choice questions (MCQs)}
\begin{enumerate}
  \item Why is pricing in e-commerce best viewed as both a prediction and an optimization problem?
  \begin{enumerate}
    \item Because prices are always random.
    \item Because the platform must estimate demand response and then choose the best feasible action under business constraints.
    \item Because forecasting and optimization are mathematically identical.
    \item Because optimization eliminates the need for data.
  \end{enumerate}

  \item Which statement best describes price elasticity?
  \begin{enumerate}
    \item It is the speed of the API that serves prices.
    \item It describes how demand responds when price changes.
    \item It is the number of SKUs in a category.
    \item It guarantees causal inference from any historical dataset.
  \end{enumerate}

  \item In promotion analytics, \emph{uplift} is most closely concerned with:
  \begin{enumerate}
    \item the additional causal effect of a promotion beyond what would have happened otherwise.
    \item the size of the ERP database.
    \item the number of products in a campaign.
    \item the compression ratio of log files.
  \end{enumerate}

  \item Which is the best example of a realistic pricing constraint?
  \begin{enumerate}
    \item The model should only use purple charts.
    \item The recommended price must preserve a minimum margin and respect contractual rules.
    \item The system should never store timestamps.
    \item The experiment should avoid all monitoring.
  \end{enumerate}

  \item Why are pricing experiments especially sensitive?
  \begin{enumerate}
    \item Because prices have no effect on customer behavior.
    \item Because pricing directly affects revenue, margin, customer trust, and sometimes legal or contractual obligations.
    \item Because they can only be run offline.
    \item Because recommendations and pricing are unrelated.
  \end{enumerate}
\end{enumerate}

\subsection*{Answer key}
\begin{enumerate}
  \item (b)
  \item (b)
  \item (a)
  \item (b)
  \item (b)
\end{enumerate}

\section{Exercises (short)}
\begin{enumerate}
  \item Compare revenue-maximizing and margin-maximizing pricing decisions for one product category. Explain when they may diverge.
  \item Design a feature set for estimating promotional uplift on a category page campaign. Include traffic, inventory, price, and seasonality signals.
  \item Propose a constrained pricing policy for a B2B account segment where negotiated price lists exist in ERP but promotional flexibility is allowed for selected categories.
\end{enumerate}

\section{Mini-case (Odoo-linked, vendor-neutral)}
\textbf{Scenario:} A hybrid B2C/B2B company uses a headless storefront, Odoo for ERP and accounting, and a data platform for forecasting. The company wants better pricing decisions for seasonal inventory and account-specific promotions.

\textbf{Task:} Design a pricing workflow:
\begin{itemize}
  \item Define which signals come from storefront behavior, which come from Odoo, and which are external.
  \item Explain how you would separate public promotional pricing from ERP-governed B2B price lists.
  \item Identify two experiment guardrails and two operational constraints that must be enforced before automated pricing recommendations are exposed.
\end{itemize}
